\section{High ramification and fixed components}
\label{app:quartic-high-ramification-fixed-components}

This appendix continues the target-span-two filtration of
\cref{app:quartic-frontier-and-ramification}.  Its conclusions are
conditional on the leading-curve and four-loci reductions stated in the
main text, and on the low-degree plane and cubic-coordinate theorems used
there.  The fixed-component argument also uses the standard fact that the
polynomial Jacobian centralizer of an indecomposable polynomial in two
variables is its polynomial ring.  Full proof notes, finite chart lists,
and exact programs are preserved in
\begin{center}
\texttt{code/program-2-2026-07-29-v2/}.
\end{center}
Every supplied program was source-reviewed and replayed successfully.  This
verifies the encoded identities and finite exceptional charts; it is not an
independent proof that the upstream geometric reductions are exhaustive.

\subsection{Primitive high ramification and edge cases}

Let
\[
 H_4=(P,Q,0),\qquad P,Q\in k[x,y]_4,\qquad \gcd(P,Q)=1,
\]
and put \(R=(H_3)_3\in k[x,y]_3\).  Define
\[
 U=J(Q,R),\qquad V=J(P,R),\qquad W=J(P,Q),
\]
and, when all three minors are nonzero,
\[
 H=\gcd(U,V,W),\qquad r=\deg H.
\]

\begin{theorem}[Primitive binary high ramification]
\label{thm:quartic-high-ramification}
Assume the Program~2 reductions place a quartic Keller map in the
primitive coprime binary-pencil locus above.  If \(r=4\) or \(r=5\), then
the map is a polynomial automorphism.
\end{theorem}

\begin{proof}
Write \(U=Hu,V=Hv,W=Hw\).  The planar minors satisfy
\[
 u\,dP-v\,dQ+w\,dR=0,                       \tag{A.1}
\]
with
\[
 \deg u=\deg v=5-r,\qquad \deg w=6-r.
\]
For \(r=5\), \(u,v\) are constants.  If \(S=uP-vQ\), then
\[
 dS=-w\,dR,\qquad dw\wedge dR=0,
\]
so homogeneity gives
\[
 R=cw^3,\qquad S=-\frac{3c}{4}w^4.
\]
The resulting aligned cube--fourth-power pair is automorphic by the
cubic-coordinate and plane-degree reduction.

For \(r=4\), first suppose \(u,v\) are dependent.  After a target change,
\[
 \ell\,dS+w\,dR=0,\qquad 4\ell S+3wR=0.     \tag{A.2}
\]
If \(\ell\mid w\), equation (A.2) again produces an aligned
cube--fourth-power pair.  If \(\ell\nmid w\), it forces
\(\ell^4T=cw^3\) with \(R=\ell T\), which is incompatible with unique
factorization.

Now suppose \(u,v\) are independent and choose \(u=x,v=y\).  Euler
contraction and differentiation of (A.1) give
\[
 P=\frac{wR_x-3Rw_x}{4},\qquad
 Q=\frac{3Rw_y-wR_y}{4}.                    \tag{A.3}
\]
If the residual quadratic \(w\) is a square, (A.3) makes \(P,Q\) share a
linear factor, contrary to primitivity.  If \(w\) is reduced, the
determinant arc reduces the first transverse normal direction to two
binary linear forms \(\ell,m\).  The necessary divisibility is
\[
 \Gamma\mid
 \frac12\bigl(\mathcal A\ell^2+
 2\mathcal B\ell m+\mathcal C m^2\bigr),
 \qquad
 \mathcal A\mathcal C-\mathcal B^2=4xy\Gamma.       \tag{A.4}
\]
For squarefree \(\Gamma\), reduction modulo \(\Gamma\) gives
\(\ell=m=0\).  The repeated-root incidence has two components: one violates
\(\gcd(P,Q)=1\), and the other is a rational conic.  The exact homogeneous
parametrization checks its affine chart, its omitted point at infinity, and
the \(2+2\) divisor, where the complete first-normal kernel is two
dimensional.  In each chart the next determinant coefficient kills the
kernel.  A later coefficient then kills the remaining quadratic normal
directions.  All nonlinear terms are binary, so a target change reduces the
map to a plane Keller pair with one coordinate of degree at most four.
\end{proof}

\begin{proposition}[Fourth-power and zero-minor boundaries]
\label{prop:quartic-fourth-power-zero-minor}
Subject to the same structural inputs:
\begin{enumerate}[label=(\roman*)]
\item if a primitive coprime leading pencil contains a fourth power, then
its cubic and quadratic normal layers route it to the binary,
quadratic-source, or aligned pure-power branch, so it creates no additional
open leaf;
\item if one of \(U,V,W\) vanishes, including the separate case \(R=0\),
then the map is a polynomial automorphism.
\end{enumerate}
\end{proposition}

\begin{proof}
For (i), normalize \(H_4=(x^4,Q,0)\), \(x\nmid Q\).  The highest
determinant equation makes \(Q(1,u,v)\) and \(R(1,u,v)\) algebraically
dependent.  Their common polynomial generator has degree one or two,
routing respectively to the binary or quadratic-source branch; a constant
generator gives an aligned pure power.  If \(R=0\), the same argument
applies to \(E=(H_2)_3\).

For (ii), if \(J(Q,R)=0\) with \(Q,R\ne0\), Euler's identity gives
\[
 3R\,dQ-4Q\,dR=0,
\]
and hence \(Q=\alpha L^4,R=\beta L^3\).  The other one-minor case is
symmetric, and two vanishing minors with \(R\ne0\) contradict target span
two.  When \(R=0\), the next determinant equation is
\[
 J(P,Q)E_z=0,
\]
so \(F_3=e(x,y)+\mu z\), \(\deg e\le2\).  Straightening this coordinate
reduces the remaining two coordinates to a plane Keller map with a
coordinate of degree at most seven; the corrected low-degree plane theorem
applies.
\end{proof}

\subsection{The corrected fixed-component valuation}

Write
\[
 P=GA,\qquad Q=GB,\qquad \gcd(A,B)=1,
\]
where \(g=\deg G\), and again put \(R=(H_3)_3\).  In this primitive branch
define
\[
 t=A/B,\qquad \frac{R^4}{Q^3}=\rho(t).
\]

\begin{lemma}[Fixed-component valuation]
\label{lem:fixed-component-valuation}
Let an irreducible component \(\Gamma\) occur in \(G\) with multiplicity
\(s\) and in the reduced special fiber \(A-\xi B\) with multiplicity
\(m\).  Put
\[
 c_\xi=\operatorname{ord}_\xi\rho
 \quad(\xi\ne\infty),\qquad
 c_\infty=\operatorname{ord}_\infty\rho+3.
\]
Then
\[
 4\nu_\Gamma(R)=3s+c_\xi m,\qquad
 \sum_{\xi\in\PP^1}c_\xi=3.                 \tag{A.5}
\]
For a nonfixed component in the same fiber, equation (A.5) holds with
\(s=0\).
\end{lemma}

\begin{proof}
At a finite fiber, \(B\) is a unit along \(\Gamma\), and
\(\nu_\Gamma(t-\xi)=m\).  Taking the valuation of
\(R^4/Q^3=\rho(t)\) gives the first identity.  At infinity, \(B\) has
order \(m\), producing the added \(3\) in the definition of
\(c_\infty\).  The second identity is the degree-zero divisor identity for
\(\rho\), with that adjustment.
\end{proof}

The term \(3s\) is essential, and the numbers \(c_\xi\) need not be
nonnegative: a pole of \(\rho\) can be cancelled by \(G^3\).  Thus any
earlier provisional closure that omitted \(3s\) or imposed
\(c_\xi\ge0\) at every fixed fiber is superseded by
\cref{lem:fixed-component-valuation}.

\begin{corollary}[Exceptional valuation at an isolated basepoint]
\label{cor:quartic-fixed-component-basepoint-valuation}
Let \(p\in\PP^2\) be a common zero of \(G,A,B\), blow up \(p\), and let
\(E\) be the exceptional divisor.  Put
\[
 g=\nu_E(G),\qquad a=\nu_E(A),\qquad b=\nu_E(B),\qquad
 r=\nu_E(R).
\]
If \(a=b\) and the residue of \(A/B\) in \(k(E)\) is not identically a
zero or pole of \(\rho\), then
\[
 4r=3(g+b),\qquad\text{and consequently}\qquad 4\mid g+b.
\]
In particular, the configuration \(g=a=b=1\) with nonconstant residual
ratio \(A/B\) is impossible.
\end{corollary}

\begin{proof}
The identity \(Q=GB\) and the primitive relation
\(R^4/Q^3=\rho(A/B)\) give
\[
 4r-3(g+b)=\nu_E\!\left(\rho(A/B)\right).
\]
Under the stated residue hypothesis, \(A/B\) is a unit at the generic point
of \(E\) and its image is neither a zero nor a pole of \(\rho\).  The
right-hand side is therefore zero.  The divisibility follows because
\(\gcd(3,4)=1\); if \(g=b=1\), it would require \(4r=6\).
\end{proof}

\begin{theorem}[Genuinely nonbinary fixed components]
\label{thm:quartic-nonbinary-fixed-components}
Assume the four-loci reduction and the polynomial-centralizer fact stated
at the beginning of this appendix.  Every genuinely nonbinary primitive
fixed-component branch is closed.
\end{theorem}

\begin{proof}
If \(\deg G=2\), equations (A.5) for a fixed line and the residual
nonfixed line give incompatible congruences modulo four, unless the whole
reduced fiber is fixed; that endpoint is the already-recorded same-fiber
rank-one chart.

If \(\deg G=1\), let \(m\in\{1,2,3\}\) be the multiplicity of the fixed
line in its reduced cubic fiber.  Multiplicity one conflicts with the
valuation on the residual factor, and multiplicity two is impossible
modulo four.  For multiplicity three, the only possibilities are the
aligned automorphic case or
\[
 H_4=(x^4,xR,0),\qquad (H_3)_3=R,\qquad x\nmid R.   \tag{A.6}
\]
The binary common-generator subcase stays in the binary tree.  In the
genuinely nonbinary subcase, the Hamiltonian derivation
\(\delta=J_{y,z}(R,-)\) has kernel \(k(x)[R]\).  The first two normal
equations give, after target changes,
\[
\begin{aligned}
 A&=ax^3,&
 B&=xE+\frac a4R-\frac\gamma4x^3,\\
 C&=\mu x^2,&
 D&=xL_3+\frac a4E-\frac\nu4x^2.
\end{aligned}                                      \tag{A.7}
\]
Reassembly yields
\[
 F_1=f(x)+\ell,\qquad
 F_2=(x+a/4)F_3+p(x)+m.
\]
The homogeneous degree-five determinant equation is
\[
 4x^3\det(dx,dm,dR)-R\det(dx,d\ell,dR)=0.          \tag{A.8}
\]
Since \(x\nmid R\), both determinants in (A.8) vanish.  Nonbinary
primitivity then forces \(\ell,m\in kx\), making two differential rows
dependent and contradicting the Keller condition.
\end{proof}

\subsection{Binary fixed factors}

\begin{theorem}[Binary fixed-factor exclusion]
\label{thm:quartic-binary-fixed-factors}
Assume the Program~2 reductions reach a binary leading pencil
\[
 H_4=(P,Q,0)
\]
with a nonconstant greatest common divisor \(G\).  If
\(\deg G=1,2,\) or \(3\), then every quartic Keller map in the corresponding
stratum is a polynomial automorphism.
\end{theorem}

\begin{proof}[Computer-assisted proof]
For \(\deg G=3\), normalize \(H_4=(xG,yG,0)\).  The three multiplicity
types of \(G\) are squarefree, double-plus-simple, and a triple line.  On
the squarefree chart, the first-normal determinant is
\[
 -6^6\operatorname{Disc}(G)\operatorname{Res}(G,R)^2.
\]
The exact root-incidence sweep closes its resultant divisor.  Separate
exact sweeps of the double-line and triple-line types either kill the normal
amplitudes or give a triangular degree drop.

For \(\deg G=1\), the valuation reduction normalizes
\[
 H_4=(x^4,xB,0),\qquad x\nmid B.
\]
There are two affine orbits for the residual cubic:
\[
 B=y^3,\qquad B=y^3+x^2y.
\]
In general coordinates the first-normal determinant is
\[
 248832\,b_3^3r_3^4\Theta,
\]
where
\[
\begin{aligned}
\Theta={}&3b_1^2r_3^2-4b_1b_2r_2r_3-6b_1b_3r_1r_3
+4b_1b_3r_2^2\\
&+4b_2^2r_1r_3-4b_2b_3r_1r_2+3b_3^2r_1^2.
\end{aligned}
\]
The exact exceptional-divisor sweep treats the projective endpoints and
the resonant vertex; the latter reassembles to a plane reduction.

For \(\deg G=2\), the coprime residual quadratic pencil is a secant of the
Veronese conic and can be normalized as
\[
 P=x^2G,\qquad Q=y^2G.
\]
Writing
\[
 G=g_0x^2+g_1xy+g_2y^2,\qquad
 R=r_0x^3+r_1x^2y+r_2xy^2+r_3y^3,
\]
the \(8\)-by-\(8\) first-normal determinant factors exactly as
\begin{equation}
\begin{aligned}
41472&(4g_0g_2-g_1^2)
(-4g_0r_1+3g_1r_0)\\
&\quad\cdot(-3g_1r_3+4g_2r_2)
\operatorname{Res}(G,R)^2.
\end{aligned}
\tag{A.9}
\end{equation}
The four divisors in (A.9), all their intersections, rank jumps,
projective endpoints, and the final aligned resonance are covered by 38
exact replay groups.  No continuity argument across an inverted factor is
used.  Off the divisors the normal matrix is invertible; on every divisor
the determinant arc either produces a unit obstruction or reduces to a
plane automorphism.
\end{proof}

\begin{corollary}[Revised quartic frontier]
\label{cor:quartic-revised-frontier}
Conditional on the complete upstream routing into the recorded Program~2
charts, the only unresolved degree-four calculation is the exceptional
triple-ramification \(F_4\) compatibility system.  A separate global
case-tree audit must still verify that the leading-curve and four-loci
reductions route every quartic map into exactly one recorded terminal
chart.
\end{corollary}

This corollary is not a proof that ordinary degree four is impossible and
does not establish \(D_{\min}\ge5\).  Independent specialist proof review
and reproduction of the new exact calculations also remain publication
requirements.
